
\documentclass[xcolor=dvipsnames, aspectratio = 169]{beamer}

\usetheme[progressbar=frametitle, numbering = fraction]{metropolis}

%% Theme Setup:
% Use Metropolis Theme, show progressbar for each section, show slide pages as fractions
\usetheme[progressbar=frametitle, numbering = fraction]{metropolis}

% Packages
\usepackage{appendixnumberbeamer}
\usepackage{pdfpages}
\usepackage{booktabs}
\usepackage{bookmark}
\usepackage[english]{babel}
\usepackage{caption}
\usepackage{amsmath}
\usepackage{xcolor}
\usepackage{amsmath}
\usepackage{comment}
\usepackage{enumitem} % to use a b c for 
\setitemize{label=$\bullet$} %enumerate
\usepackage[export]{adjustbox} % to use left in includegraphics
\usepackage{fixmath} % to write math expressions in bold
\usepackage{xspace}
\usepackage[scale=2]{ccicons}
\usepackage{csquotes}
\usepackage{subcaption} % needed for making subfigures
\usepackage{pifont} % ty
%Needed for accessing dinglist for making check marks and crosses
\usepackage[round]{natbib} % important for using citet and citep !!!
\usepackage{bbm} % for printing the indicator sign properly
\usepackage{dsfont} % too
\usepackage{verbatim}
%%%% For Making Metropolis nice %%%%%
\setbeamertemplate{mini frame}{}
\setbeamertemplate{mini frame in current section}{}
\setbeamertemplate{mini frame in current subsection}{}

% To make sections on top of each slide but overwrite dots \makeatletter
\setbeamertemplate{headline}{%
\begin{beamercolorbox}[sep=0.1cm, ht=2.3em]{section in head/foot}
		\vskip2pt\insertnavigation{\paperwidth}\vskip-2pt
\end{beamercolorbox}%
\begin{beamercolorbox}[colsep=1.5pt]{lower separation line head}
\end{beamercolorbox}}

% To make box for titles of each slide
\setbeamertemplate{frametitle}{
	\nointerlineskip
	\begin{beamercolorbox}[sep=0.2cm, ht=1.8em,wd=\paperwidth]{frametitle}
		\vbox{}\vskip-1ex%
		\strut\insertframetitle\strut
		\vskip-0.5ex%
	\end{beamercolorbox}
	\usebeamertemplate*{progress bar in head/foot}}

% Change the color of box at top (fg: text, bg, box color)
% Lower Box with title
\definecolor{black2}{RGB}{27,55,61}
\setbeamercolor{section in head/foot}{fg=White, bg=black2}  % Upper Box with sections

%% For alert block
\setbeamercolor{block title alerted}{fg= White, bg=black2}
%body
\setbeamercolor{block body alerted}{fg=black!90,bg=lightgray}


% Change font size of title of each slide 
\setbeamerfont{frametitle}{size=\small,series=\bfseries}
\setbeamerfont{section in head/foot}{size=\tiny}

% This section adds the navigation dots for all slides
\usepackage{remreset}% tiny package containing just the \@removefromreset command
\makeatletter
\@removefromreset{subsection}{section}
\beamer@compresstrue					% Locates dots nicely below section names
\makeatother
\setcounter{subsection}{1}

% Needed for metropolis
\newcommand{\themename}{\textbf{\textsc{metropolis}}\xspace}
\newcommand*{\captionsource}[2]{\caption[{#1}]{#1\hspace{\linewidth}\textbf{Source:
		}#2}}

%% Further Specifications %%

% To not print "References" at the slide where references are
\renewcommand{\bibsection}{}

% To make \widebar possible without using mathabx package (produces further problems)
\makeatletter
\newcommand*\rel@kern[1]{\kern#1\dimexpr\macc@kerna}
\newcommand*\widebar[1]{%
	\begingroup
	\def\mathaccent##1##2{%
		\rel@kern{0.8}%
		\overline{\rel@kern{-0.8}\macc@nucleus\rel@kern{0.2}}%
		\rel@kern{-0.2}%
	}%
	\macc@depth\@ne
	\let\math@bgroup\@empty \let\math@egroup\macc@set@skewchar
	\mathsurround\z@ \frozen@everymath{\mathgroup\macc@group\relax}%
	\macc@set@skewchar\relax
	\let\mathaccentV\macc@nested@a
	\macc@nested@a\relax111{#1}%
	\endgroup
}
\makeatother
%%%%%%%%%%%%%%%%%%%%%%

% Set Graphics path
\graphicspath{{figures/}}

\title{Optimal flood policy mix}

\date{March 26, 2026}
\author{Taco Prins, Lennard Schlattmann, Daniel Schmidt}

%\institute{DNB internal}

\begin{document}
\maketitle

{ % to make sections appear in navigation bar at top
\AtBeginSection{} % but not as section slides, 
% close at the end

\section{Motivation}

\begin{frame}{Motivation}
    \begin{itemize}
    	\setlength{\itemsep}{0.3cm}
        \item Households have \textbf{heterogeneous flood risk beliefs}, some underestimate risk
        \item Mandatory flood risk insurance often legally/politically infeasible 
        \item Several \textbf{policy alternatives:} 
        \begin{enumerate}[label=\arabic*.] 
            \item Subsidies for insurance correct under-insurance
            \item Subsidies for private adaptation measures correct under-investment in adaptation
            \item Disaster relief programs provide safety net for uninsured households, but reduce incentives to insure or adapt
        \end{enumerate}
        \item \textbf{Questions:} What is the optimal policy mix? How does it depend on distribution of flood risk beliefs, risk aversion, other parameters?
        \item \textbf{Approach:} Stylized theoretical model with sufficient statistics approach
    \end{itemize}
\end{frame}

\section{Preliminary results}

\begin{frame}{Preliminary results}

    Focus on trade-off between insurance subsidy and disaster relief:
    \begin{align}
    \text{MVPF}_s 
    &= \frac{
    \overbrace{(p-q^*)\,\Delta u \cdot \dfrac{\partial I}{\partial s}}^{\text{welfare gain from take-up}}
    \;+\;
    \overbrace{pd\,u'(c_{ins})\cdot I}^{\text{windfall to insured}}
    }{
    \underbrace{pd\!\left[I + (s-a)\,\dfrac{\partial I}{\partial s}\right]}_{\text{net fiscal cost}}
    }
    \label{eq:mvpf_s}
    \\[14pt]
    \text{MVPF}_a 
    &= \frac{
    \overbrace{(p-q^*)\,\Delta u \cdot \dfrac{\partial I}{\partial a}}^{\text{welfare loss from crowd-out}}
    \;+\;
    \overbrace{pd\,u'(c_f)\cdot U}^{\text{transfer to uninsured}}
    }{
    \underbrace{pd\!\left[U + (s-a)\,\dfrac{\partial I}{\partial a}\right]}_{\text{net fiscal cost}}
    }
    \label{eq:mvpf_a}
    \end{align}
    
\end{frame}

\section{Questions}

\begin{frame}{Questions}
  \begin{itemize}
    \item Is our methodological approach a good idea for this question?
    \item What should we focus on first?
    \begin{itemize}
      \item Subsidies for insurance and disaster relief, no private adaptation?
      \item Location choice or not?
    \end{itemize}
    \item Underestimation of risk
    \begin{itemize}
      \item There are several possible reasons for underestimation of risk (cognitive biases, lack of information, ...). Does it matter for our analysis what the source of underestimation is?
      \item Possible counterargument: better information about flood risk is more natural policy response to biased beliefs than subsidies or disaster relief
    \end{itemize}
  \end{itemize}
\end{frame}

% ==========================================================

\section{Backup slides}

% ==========================================================

\begin{frame}{Model: Locations, preferences, and risk}
\begin{itemize}[leftmargin=0.0cm]
    \setlength{\itemsep}{0.3cm}
    \item[] \textbf{Two locations:} inland $l$ and waterfront $w$ with housing stocks $H_l(p_l)$ and $H_w(p_w)$
\item[] \textbf{Agents:} income $y$ drawn from continuous distribution; idiosyncratic location taste $\varepsilon_{ij}$
\item[] \textbf{Utility:} $u(s,x) + \varepsilon_{ij}$, with concave $u(\cdot)$ and Inada-type conditions
\item[] \textbf{Housing services:}
\[
s =
\begin{cases}
h & \text{if } j=l \\
\omega h & \text{if } j=w \quad (\omega>1 \text{ amenity})
\end{cases}
\]
\item[] \textbf{Flood risk:} flood occurs idiosyncratically with true probability $p$, damage proportional to housing: $dh$
\item[] \textbf{Beliefs:} subjective flood probability $q\in[0,1]$ with distribution (unknown) over $(q,y)$
\end{itemize}
\end{frame}

% ==========================================================
\begin{frame}{Model: Timing and choices}
\textbf{One-period model; steady-state comparison}\\
\begin{enumerate}[label=\arabic*.] 
  \item Household observes $(y,q)$ and draws $\varepsilon_{ij}$
  \item Household chooses location $j\in\{l,w\}$ and housing $h$
  \item If $j=w$, household chooses whether to buy flood insurance
  \item Flood realization occurs (with true probability $p$)
  \item Consumption of numeraire good $x$ adjusts to satisfy the budget constraint
\end{enumerate}
\vspace{0.4em}
\textbf{Policy instruments:}
\begin{itemize}
  \item Proportional income tax $\tau$ finances policy
  \item Disaster aid generosity $a$ (reduces uncovered repair cost in flood state)
  \item Insurance subsidy rate $s$ (reduces insurance premium paid by household)
\end{itemize}
\end{frame}

% ==========================================================
\begin{frame}{Model: Budget constraints by choice}
Let $p_l$ and $p_w$ denote (long-run) housing prices inland and waterfront.

\medskip
\textbf{(i) Inland:}
\[
p_l h_i^l + x_i = (1-\tau) y_i
\]

\textbf{(ii) Waterfront with insurance (actuarially fair premium):}
\[
\qquad p_w h_i^w + (1-s)p d h_i^w + x_i = (1-\tau) y_i, \quad \text{independent of flood realization}
\]

\textbf{(iii) Waterfront without insurance:}
\[
\begin{aligned}
\text{No flood: } & p_w h_i^w + x_i = (1-\tau)y_i,\\
\text{Flood: } & p_w h_i^w + (1-a)d h_i^w + x_i = (1-\tau) y_i
\end{aligned}
\]
\end{frame}

% ==========================================================
\begin{frame}{Model: Indirect utilities and belief-based decision making}
Define indirect utility conditional on location/insurance choice:

\begin{itemize}
  \item Inland: $V^l(y; p_l)$.
  \item Waterfront insured: $V^{w,\text{yes}}(y; p_w)$.
  \item Waterfront uninsured: households maximize \emph{subjective} expected utility using $q$:
  \[
  \tilde V^{w,\text{no}}(q,y; p_w) = \max_{h^w} (1-q)u(\omega h^w, x_0) + q\,u(\omega h^w, x_1),
  \]
  with $x_0=(1-\tau)y - p_w h^w$ and $x_1=(1-\tau)y-(p_w+(1-a)d)h^w$.
\end{itemize}

\medskip
\textbf{Key friction:} planner evaluates outcomes with \emph{true} $p$, but households choose $h^w$ and insurance based on $q$. This creates a wedge between perceived optimal and socially optimal welfare

\end{frame}

% ==========================================================
\begin{frame}{Findings: Insurance take-up threshold and comparative statics}
There exists a \textbf{belief threshold} $q^*(y)\in(0,1)$ such that:
\[
\text{Buy insurance} \iff q \ge q^*(y)
\]

\smallskip
\textbf{Intuition for threshold:}
\begin{itemize}
  \item If $q=0$: insurance is never purchased (costly, no perceived benefit)
  \item If perceived expected loss is large enough, concavity implies insurance is preferred
  \item Disaster aid $a$ makes being uninsured less painful in the flood state $\rightarrow$ \textbf{reduces insurance incentives}
  \item Insurance subsidy $s$ reduces the premium $\rightarrow$ \textbf{raises insurance incentives}
\end{itemize}

\medskip
\textbf{Comparative statics (qualitative):}
\[
\frac{\partial q^*(y)}{\partial a} > 0,
\qquad
\frac{\partial q^*(y)}{\partial s} < 0
\]

\end{frame}

% ==========================================================
\begin{frame}{Findings: Location choice}
Type-1 extreme value tastes $(\varepsilon_{ij})$ imply logit structure for waterfront choice

\medskip
\textbf{Uninsured waterfront choice probability (for $q<q^*(y)$):}
\[
P^{w,\text{no}}(q,y)=
\frac{\exp\left(\tilde V^{w,\text{no}}(q,y)\right)}
{\exp\left(\tilde V^{w,\text{no}}(q,y)\right)+\exp\left(V^l(y)\right)}.
\]

\textbf{Insured waterfront choice probability (for $q\ge q^*(y)$):}
\[
P^{w,\text{yes}}(y)=
\frac{\exp\left(V^{w,\text{yes}}(y)\right)}
{\exp\left(V^{w,\text{yes}}(y)\right)+\exp\left(V^l(y)\right)}.
\]

\medskip
\textbf{Implication:} Policy affects welfare through two channels:
\begin{itemize}
  \item \textbf{Intensive margin:} housing demand $h$ and insurance choice
  \item \textbf{Extensive margin:} shifts in location probability into the waterfront area
\end{itemize}
\end{frame}

% ==========================================================
\begin{comment}
\begin{frame}{First findings: Social planner: evaluates welfare under true flood probability $p$ (not $q$)}
\small
\textbf{Key objects:}
\begin{itemize}
  \item Joint density of beliefs and income: $f(q,y)$
  \item Insurance threshold: $q^*(y)$, with mass of insured types $m(y)=\int_{q^*(y)}^{1} f(q,y)\,dq$
  \item Choice probabilities: $P^{w,\text{no}}(q,y)$ (uninsured, $q<q^*(y)$), $P^{w,\text{yes}}(y)$ (insured, $q\ge q^*(y)$)
\end{itemize}  

\textbf{Welfare components:}
\begin{enumerate}[label=\arabic*.] 
  \item \textbf{Expected indirect utility} for types with $q<q^*(y)$:
  \[
  \tilde V^{\text{no}}(q,y)=\log\!\Big[\exp\big(\tilde V^{w,\text{no}}(q,y)\big)+\exp\big(V^l(y)\big)\Big]
  \]
  \item \textbf{Expected indirect utility} for types with $q\ge q^*(y)$:
  \[
  V^{\text{yes}}(y)=\log\!\Big[\exp\big(V^{w,\text{yes}}(y)\big)+\exp\big(V^l(y)\big)\Big]
  \]
  \item \textbf{Belief-correction term} for uninsured waterfront movers:
  \[
  \Delta(q,y)=V^{w,\text{no}}(q,y)-\tilde V^{w,\text{no}}(q,y)
  = (p-q)\Big[u(\omega h^{w,\text{no}},x_1)-u(\omega h^{w,\text{no}},x_0)\Big]
  \]
\end{enumerate}

\smallskip

\end{frame}
\end{comment}
% ==========================================================
\begin{frame}{First findings: Social planner}
\small
\textbf{Planner objective (schematic):}
\begin{eqnarray*}
    W(a,s)&=&
\underbrace{\int_y \int_0^{q^*(y)} \tilde V^{\text{no}}(q,y) f\,dq\,dy}_{\text{uninsured (by $q$)}}
+\underbrace{\int_y V^{\text{yes}}(y) m(y)\,dy}_{\text{insured}} \\
&+&\underbrace{\int_y \int_0^{q^*(y)} P^{w,\text{no}}(q,y)\,\Delta(q,y)\, f\,dq\,dy}_{\text{belief correction}}
-\lambda_B\, \underbrace{G(a,s)}_{\text{balanced budget}} 
\end{eqnarray*}


\textbf{Balanced budget:}
\[
\int_y\!\!\int_q \tau y\, f\,dq\,dy
=
\underbrace{p\,a\,d\!\int_y\!\!\int_0^{q^*(y)} \! P^{w,\text{no}}(q,y)\,h^{w,\text{no}}\, f\,dq\,dy}_{\text{disaster-aid outlays}}
+
\underbrace{p\,s\,d\!\int_y \! P^{w,\text{yes}}(y)\,h^{w,\text{yes}}\, m(y)\,dy}_{\text{insurance-subsidy outlays}}
\]
\end{frame}


\begin{frame}{First findings: Policy effects (qualitative):}
\small
\begin{itemize}[leftmargin=-0.2cm]
  \item \textbf{Insurance subsidy $s$:}
  \begin{itemize}
    \item Lowers net premium $\rightarrow$ raises $V^{w,\text{yes}}$, $P^{w,\text{yes}}$ and $h^{w,\text{yes}}$, lowers threshold $q^*(y)$
    \item Budget cost scales with $p\,s\,d$ and insured waterfront exposure
  \end{itemize}
  \item \textbf{Disaster aid $a$:}
  \begin{itemize}
    \item Softens uninsured flood state $\rightarrow$ raises $\tilde V^{w,\text{no}}$, $P^{w,\text{no}}$, and threshold $q^*(y)$ 
    \item \textbf{Belief friction cost:} shifts marginal agents into risky area without fully internalizing true $p$
  \end{itemize}
\end{itemize}
\smallskip
\textbf{Optimal mix intuition:}
\begin{enumerate}[label=\arabic*.] 
  \item Use $s$ to correct under-insurance driven by biased beliefs \& risk aversion (ex ante)
  \item Limit $a$ to avoid excessive sorting into risky areas (moral hazard), while preserving social insurance motives (ex post)
  \item Respect the government budget and weigh the belief-correction gains against fiscal costs and induced location shifts
\end{enumerate}
\end{frame}

% ----------------------------------------------------------
\begin{frame}{Next steps:}
\textbf{Structural:}
\begin{itemize}
  \item .
\end{itemize}

\textbf{Empirical:}
\begin{itemize}
  \item .
\end{itemize}

\end{frame}

\begin{frame}{Literature}
    \begin{itemize}
        \item Flood risks: \citet{bakkensen2022}
        \item Flood insurance: \citet{hennighausen2023}, \citet{collier2023}
    \end{itemize}
\end{frame}

\appendix

\section{Appendix}

%%% REFERENCES %%%

\setbeamertemplate{frametitle continuation}{}
% to get rid of numeration in references frame title:
% i, ii, etc.

\begin{frame}[allowframebreaks]{References}
	\small
	\bibliography{subfiles/references.bib}
	\bibliographystyle{apalike}
\end{frame}

}



% ==========================================================


\end{document}
